% ============================================================================
% APPENDIX XI: FORMAL-MODELS - Mathematical Models of Complementarity
% ============================================================================
% Category: FORMAL (Pure Mathematics)
% Language: English
% Version: 1.0
% Status: FOUNDATIONAL
% Dependencies: X (FORMAL-FOUND)
% Domain: Universal (domain-independent model catalog)
% ============================================================================

\chapter{Mathematical Models of Complementarity: A Comprehensive Catalog}
\label{app:formal-models}

\begin{abstract}
This appendix presents the key concepts and formalizations for Mathematical Models of Complementarity: A Comprehensive Catalog.
It provides theoretical foundations, practical applications, and integration
points with other components of the Evidence-Based Framework (EBF).
The content supports decision-making and behavioral analysis in the context
of complementarity and context-dependent outcomes.
\end{abstract}

\begin{tcolorbox}[colback=blue!5!white, colframe=blue!75!black,
    title=Quick Reference: Mathematical Models of Complementarity: A Comprehensive Catalog]

\textbf{Appendix:} MDL (FORMAL)

\textbf{Core Concepts:}
\begin{itemize}[nosep]
\item Complementarity ($\gamma > 0$): Components interact synergistically
\item Context ($\Psi$): Situational factors that influence behavior
\item Utility dimensions: FEPSDE (Financial, Emotional, Physical, Social, Development, Experiential)
\end{itemize}

\textbf{Key Formulas:}
\begin{equation}
U = \sum_d \omega_d \cdot u_d + \gamma \cdot f(\text{interactions})
\end{equation}

\textbf{Cross-References:} Appendix UNMAPPED_B (CORE-HOW), V (CORE-WHEN), G (Glossary)
\end{tcolorbox}





% ----------------------------------------------------------------------------
% HEADER BLOCK
% ----------------------------------------------------------------------------
\begin{tcolorbox}[colback=blue!5!white,colframe=blue!75!black,title=\textbf{Appendix UNMAPPED_MDL: FORMAL-MODELS}]
\textbf{Full Title:} Mathematical Models of Complementarity: A Comprehensive Catalog

\textbf{Category:} FORMAL (Pure Mathematics - Layer 2)

\textbf{Purpose:} Catalogs all major mathematical models for representing complementarity, their assumptions, domains of validity, and inter-relationships.

\textbf{Prerequisite:} Appendix UNMAPPED_FND (FORMAL-FOUND) establishes WHAT complementarity is. This appendix shows HOW to model it mathematically.

\textbf{Key Principle:} Different models are appropriate for different mathematical contexts. No single model is universally best.
\end{tcolorbox}

% ----------------------------------------------------------------------------
% MODEL OVERVIEW
% ----------------------------------------------------------------------------
\begin{tcolorbox}[colback=gray!5!white,colframe=gray!75!black,title=\textbf{The Six Canonical Models}]
\small
\begin{enumerate}[noitemsep]
    \item \textbf{Model A:} Supermodularity (Lattice-Based)
    \item \textbf{Model B:} Cross-Partial Derivatives (Calculus-Based)
    \item \textbf{Model C:} Factorial Interaction (Statistical)
    \item \textbf{Model D:} Production Function (Multiplicative)
    \item \textbf{Model E:} Graph-Theoretic (Network)
    \item \textbf{Model F:} Tensor Representation (Multilinear)
\end{enumerate}

Each model has: Definition, Assumptions, Domain, Properties, Limitations.
\end{tcolorbox}

% ============================================================================
% MODEL A: SUPERMODULARITY
% ============================================================================
\section{Model A: Supermodularity (Lattice-Based)}
\label{sec:model-a}

\begin{tcolorbox}[colback=green!5!white,colframe=green!75!black,title=\textbf{Model A Overview}]
\textbf{Core Idea:} Complementarity as a property of functions on lattices.

\textbf{Mathematical Setting:} Lattice $(L, \leq)$ with function $f: L \to \mathbb{R}$.

\textbf{Complementarity Definition:}
\begin{equation}
f(x \vee y) + f(x \wedge y) \geq f(x) + f(y) \quad \text{(supermodular)}
\end{equation}

\textbf{Best For:} Discrete optimization, comparative statics, game theory.
\end{tcolorbox}

\subsection{Formal Definition}

\begin{definition}[Model A: Supermodular Complementarity]
Let $(L, \leq)$ be a lattice and $f: L \to \mathbb{R}$. The complementarity between elements $x$ and $y$ is:
\begin{equation}
\gamma^A_{xy} = f(x \vee y) + f(x \wedge y) - f(x) - f(y)
\label{eq:model-a}
\end{equation}
$f$ exhibits \textbf{global complementarity} if $\gamma^A_{xy} \geq 0$ for all $x, y \in L$.
\end{definition}

\subsection{Key Properties}

\begin{theorem}[Properties of Model A]
For supermodular $f$:
\begin{enumerate}
    \item \textbf{Closure under addition:} If $f, g$ supermodular, so is $f + g$
    \item \textbf{Closure under limits:} Pointwise limit of supermodular functions is supermodular
    \item \textbf{Optimization structure:} $\arg\max f$ forms a sublattice
    \item \textbf{Monotone comparative statics:} Optimal solutions increase with parameters
\end{enumerate}
\end{theorem}

\subsection{Assumptions and Limitations}

\textbf{Required Assumptions:}
\begin{itemize}
    \item Domain must be a lattice (partial order with join/meet)
    \item Complementarity is global (all pairs, not just some)
\end{itemize}

\textbf{Limitations:}
\begin{itemize}
    \item Cannot represent mixed complementarity/substitution
    \item Requires lattice structure (not always natural)
    \item Binary: either supermodular or not
\end{itemize}

% ============================================================================
% MODEL B: CROSS-PARTIAL DERIVATIVES
% ============================================================================
\section{Model B: Cross-Partial Derivatives (Calculus-Based)}
\label{sec:model-b}

\begin{tcolorbox}[colback=green!5!white,colframe=green!75!black,title=\textbf{Model B Overview}]
\textbf{Core Idea:} Complementarity as the sign of the second cross-partial derivative.

\textbf{Mathematical Setting:} Twice-differentiable function $f: \mathbb{R}^n \to \mathbb{R}$.

\textbf{Complementarity Definition:}
\begin{equation}
\gamma^B_{ij} = \frac{\partial^2 f}{\partial x_i \partial x_j}
\end{equation}

\textbf{Best For:} Continuous optimization, economic models, physics.
\end{tcolorbox}

\subsection{Formal Definition}

\begin{definition}[Model B: Differential Complementarity]
For $f: \mathbb{R}^n \to \mathbb{R}$ with $f \in C^2$, the complementarity between dimensions $i$ and $j$ is:
\begin{equation}
\gamma^B_{ij}(\mathbf{x}) = \frac{\partial^2 f}{\partial x_i \partial x_j}\bigg|_{\mathbf{x}}
\label{eq:model-b}
\end{equation}
\begin{itemize}
    \item $\gamma^B_{ij} > 0$: Complements at $\mathbf{x}$
    \item $\gamma^B_{ij} < 0$: Substitutes at $\mathbf{x}$
    \item $\gamma^B_{ij} = 0$: Independent at $\mathbf{x}$
\end{itemize}
\end{definition}

\subsection{Matrix Representation: The Hessian}

\begin{definition}[Complementarity Matrix]
The \textbf{Hessian} $\mathbf{H}$ captures all pairwise complementarities:
\begin{equation}
\mathbf{H}(\mathbf{x}) = \begin{pmatrix}
\frac{\partial^2 f}{\partial x_1^2} & \gamma^B_{12} & \cdots & \gamma^B_{1n} \\
\gamma^B_{21} & \frac{\partial^2 f}{\partial x_2^2} & \cdots & \gamma^B_{2n} \\
\vdots & \vdots & \ddots & \vdots \\
\gamma^B_{n1} & \gamma^B_{n2} & \cdots & \frac{\partial^2 f}{\partial x_n^2}
\end{pmatrix}
\end{equation}
Off-diagonal elements are complementarities; diagonal elements are curvatures.
\end{definition}

\subsection{Key Properties}

\begin{theorem}[Properties of Model B]
\begin{enumerate}
    \item \textbf{Symmetry:} $\gamma^B_{ij} = \gamma^B_{ji}$ (Young's theorem)
    \item \textbf{Local:} $\gamma^B_{ij}(\mathbf{x})$ can vary with $\mathbf{x}$
    \item \textbf{Sign can change:} Complements at one point, substitutes at another
    \item \textbf{Link to Model A:} If $\gamma^B_{ij} \geq 0$ everywhere, $f$ is supermodular
\end{enumerate}
\end{theorem}

\subsection{Assumptions and Limitations}

\textbf{Required Assumptions:}
\begin{itemize}
    \item Function must be twice continuously differentiable ($C^2$)
    \item Domain is open subset of $\mathbb{R}^n$
\end{itemize}

\textbf{Limitations:}
\begin{itemize}
    \item Not defined for discrete domains
    \item Local definition may miss global patterns
    \item Requires smoothness assumption
\end{itemize}

% ============================================================================
% MODEL C: FACTORIAL INTERACTION
% ============================================================================
\section{Model C: Factorial Interaction (Statistical)}
\label{sec:model-c}

\begin{tcolorbox}[colback=green!5!white,colframe=green!75!black,title=\textbf{Model C Overview}]
\textbf{Core Idea:} Complementarity as interaction effect in ANOVA.

\textbf{Mathematical Setting:} $2^k$ factorial design with outcome $Y$.

\textbf{Complementarity Definition:}
\begin{equation}
\gamma^C_{AB} = \bar{Y}_{11} - \bar{Y}_{10} - \bar{Y}_{01} + \bar{Y}_{00}
\end{equation}

\textbf{Best For:} Experimental design, causal inference, empirical estimation.
\end{tcolorbox}

\subsection{Formal Definition}

\begin{definition}[Model C: Factorial Complementarity]
In a $2 \times 2$ factorial design with factors $A, B \in \{0, 1\}$:
\begin{equation}
\gamma^C_{AB} = E[Y | A=1, B=1] - E[Y | A=1, B=0] - E[Y | A=0, B=1] + E[Y | A=0, B=0]
\label{eq:model-c}
\end{equation}
\end{definition}

\subsection{Linear Model Representation}

\begin{definition}[ANOVA Model]
\begin{equation}
Y = \mu + \alpha_A \cdot A + \alpha_B \cdot B + \gamma^C_{AB} \cdot (A \times B) + \epsilon
\end{equation}
where:
\begin{itemize}
    \item $\mu$: Grand mean
    \item $\alpha_A, \alpha_B$: Main effects
    \item $\gamma^C_{AB}$: Interaction effect (complementarity)
    \item $\epsilon$: Error term
\end{itemize}
\end{definition}

\subsection{Higher-Order Extensions}

For $k$ factors:
\begin{equation}
Y = \mu + \sum_i \alpha_i X_i + \sum_{i<j} \gamma_{ij} X_i X_j + \sum_{i<j<k} \gamma_{ijk} X_i X_j X_k + \cdots
\end{equation}

\begin{definition}[Three-Way Interaction]
\begin{equation}
\gamma^C_{ABC} = (\gamma_{AB|C=1} - \gamma_{AB|C=0})
\end{equation}
"Does the complementarity between $A$ and $B$ depend on $C$?"
\end{definition}

\subsection{Assumptions and Limitations}

\textbf{Required Assumptions:}
\begin{itemize}
    \item Factors are discrete (typically binary)
    \item Randomization or conditional independence
    \item Linear additivity in main effects
\end{itemize}

\textbf{Limitations:}
\begin{itemize}
    \item Does not capture continuous variation
    \item Combinatorial explosion with many factors
    \item Assumes linearity
\end{itemize}

% ============================================================================
% MODEL D: PRODUCTION FUNCTION
% ============================================================================
\section{Model D: Production Function (Multiplicative)}
\label{sec:model-d}

\begin{tcolorbox}[colback=green!5!white,colframe=green!75!black,title=\textbf{Model D Overview}]
\textbf{Core Idea:} Complementarity embedded in functional form.

\textbf{Mathematical Setting:} Production/utility functions with interaction terms.

\textbf{Complementarity Definition:} Parameter in functional form.

\textbf{Best For:} Economic modeling, structural estimation.
\end{tcolorbox}

\subsection{Cobb-Douglas Form}

\begin{definition}[Cobb-Douglas]
\begin{equation}
f(x_1, x_2) = A x_1^{\alpha} x_2^{\beta}
\end{equation}
Complementarity: $\gamma^D_{12} = \frac{\partial^2 f}{\partial x_1 \partial x_2} = A \alpha \beta x_1^{\alpha-1} x_2^{\beta-1}$

Note: Always positive if $\alpha, \beta > 0$—implies universal complementarity.
\end{definition}

\subsection{CES Form}

\begin{definition}[Constant Elasticity of Substitution]
\begin{equation}
f(x_1, x_2) = \left( \alpha x_1^\rho + (1-\alpha) x_2^\rho \right)^{1/\rho}
\end{equation}
\begin{itemize}
    \item $\rho > 1$: Complements
    \item $\rho = 1$: Linear (perfect substitutes)
    \item $\rho < 1$: Gross substitutes
    \item $\rho \to 0$: Cobb-Douglas
    \item $\rho \to -\infty$: Leontief (perfect complements)
\end{itemize}
\end{definition}

\subsection{Translog Form}

\begin{definition}[Translog (Transcendental Logarithmic)]
\begin{equation}
\ln f = \alpha_0 + \sum_i \alpha_i \ln x_i + \frac{1}{2} \sum_i \sum_j \gamma^D_{ij} \ln x_i \ln x_j
\label{eq:translog}
\end{equation}
Here $\gamma^D_{ij}$ is directly estimated as a parameter.
\end{definition}

\begin{theorem}[Translog Complementarity]
In the translog model:
\begin{equation}
\frac{\partial^2 \ln f}{\partial \ln x_i \partial \ln x_j} = \gamma^D_{ij}
\end{equation}
$\gamma^D_{ij}$ is the elasticity of complementarity.
\end{theorem}

\subsection{Assumptions and Limitations}

\textbf{Required Assumptions:}
\begin{itemize}
    \item Specific functional form (parametric)
    \item Inputs are continuous and positive
    \item Function is monotone increasing
\end{itemize}

\textbf{Limitations:}
\begin{itemize}
    \item Results depend on chosen functional form
    \item May impose complementarity/substitution by construction
    \item Global restrictions may not hold locally
\end{itemize}

% ============================================================================
% MODEL E: GRAPH-THEORETIC
% ============================================================================
\section{Model E: Graph-Theoretic (Network)}
\label{sec:model-e}

\begin{tcolorbox}[colback=green!5!white,colframe=green!75!black,title=\textbf{Model E Overview}]
\textbf{Core Idea:} Complementarity as weighted edges in a graph.

\textbf{Mathematical Setting:} Graph $G = (V, E)$ with edge weights.

\textbf{Complementarity Definition:} Edge weight $\gamma_{ij} = w(i, j)$.

\textbf{Best For:} Networks, social interactions, spatial models.
\end{tcolorbox}

\subsection{Formal Definition}

\begin{definition}[Model E: Graph Complementarity]
Given a weighted graph $G = (V, E, w)$:
\begin{itemize}
    \item Nodes $V = \{1, \ldots, n\}$ are elements
    \item Edge $(i, j) \in E$ exists if $i$ and $j$ interact
    \item Weight $\gamma^E_{ij} = w(i, j)$ is the complementarity strength
\end{itemize}
\end{definition}

\subsection{Adjacency Matrix Representation}

\begin{definition}[Complementarity Adjacency Matrix]
\begin{equation}
\mathbf{W} = \begin{pmatrix}
0 & \gamma^E_{12} & \cdots & \gamma^E_{1n} \\
\gamma^E_{21} & 0 & \cdots & \gamma^E_{2n} \\
\vdots & \vdots & \ddots & \vdots \\
\gamma^E_{n1} & \gamma^E_{n2} & \cdots & 0
\end{pmatrix}
\end{equation}
\end{definition}

\subsection{Graph Laplacian}

\begin{definition}[Complementarity Laplacian]
\begin{equation}
\mathbf{L} = \mathbf{D} - \mathbf{W}
\end{equation}
where $\mathbf{D} = \text{diag}(\sum_j \gamma_{1j}, \ldots, \sum_j \gamma_{nj})$.
\end{definition}

\begin{theorem}[Spectral Properties]
The eigenvalues of $\mathbf{L}$ characterize complementarity structure:
\begin{itemize}
    \item Smallest eigenvalue is 0 (connected components)
    \item Second smallest (Fiedler value) measures connectivity
    \item Eigenvectors identify complementarity clusters
\end{itemize}
\end{theorem}

\subsection{Assumptions and Limitations}

\textbf{Required Assumptions:}
\begin{itemize}
    \item Pairwise interactions only (no higher-order without hypergraphs)
    \item Fixed network structure
\end{itemize}

\textbf{Limitations:}
\begin{itemize}
    \item Does not naturally capture functional dependence
    \item Network must be specified exogenously
\end{itemize}

% ============================================================================
% MODEL F: TENSOR REPRESENTATION
% ============================================================================
\section{Model F: Tensor Representation (Multilinear)}
\label{sec:model-f}

\begin{tcolorbox}[colback=green!5!white,colframe=green!75!black,title=\textbf{Model F Overview}]
\textbf{Core Idea:} Complementarity as a multi-dimensional tensor.

\textbf{Mathematical Setting:} Tensor $\boldsymbol{\Gamma}$ of arbitrary order.

\textbf{Complementarity Definition:} Tensor entries $\Gamma_{i_1 i_2 \cdots i_k}$.

\textbf{Best For:} Higher-order interactions, multi-dimensional data.
\end{tcolorbox}

\subsection{Formal Definition}

\begin{definition}[Model F: Tensor Complementarity]
A \textbf{complementarity tensor} of order $k$ is:
\begin{equation}
\boldsymbol{\Gamma} \in \mathbb{R}^{n \times n \times \cdots \times n} \quad (k \text{ indices})
\end{equation}
where $\Gamma_{i_1 i_2 \cdots i_k}$ is the $k$-way complementarity among elements $i_1, \ldots, i_k$.
\end{definition}

\subsection{Order Hierarchy}

\begin{itemize}
    \item \textbf{Order 2:} Matrix $\Gamma_{ij}$ (pairwise complementarity)
    \item \textbf{Order 3:} 3-tensor $\Gamma_{ijk}$ (three-way complementarity)
    \item \textbf{Order $k$:} $k$-tensor (general higher-order)
\end{itemize}

\subsection{Tensor Decomposition}

\begin{theorem}[CP Decomposition]
Any complementarity tensor can be decomposed as:
\begin{equation}
\Gamma_{i_1 \cdots i_k} = \sum_{r=1}^{R} \lambda_r \cdot a^{(1)}_{i_1 r} \cdot a^{(2)}_{i_2 r} \cdots a^{(k)}_{i_k r}
\end{equation}
where $R$ is the tensor rank and $\mathbf{a}^{(j)}$ are factor matrices.
\end{theorem}

This allows dimension reduction for high-order complementarity.

\subsection{Symmetry Constraints}

\begin{definition}[Symmetric Tensor]
A complementarity tensor is \textbf{symmetric} if:
\begin{equation}
\Gamma_{i_{\pi(1)} \cdots i_{\pi(k)}} = \Gamma_{i_1 \cdots i_k}
\end{equation}
for all permutations $\pi$.
\end{definition}

\subsection{Assumptions and Limitations}

\textbf{Required Assumptions:}
\begin{itemize}
    \item Fixed set of elements
    \item Defined interaction order
\end{itemize}

\textbf{Limitations:}
\begin{itemize}
    \item Computational complexity grows exponentially with order
    \item Estimation requires large samples
    \item Interpretation becomes difficult for high order
\end{itemize}

% ============================================================================
% MODEL COMPARISON
% ============================================================================
\section{Model Comparison and Equivalences}
\label{sec:comparison}

\subsection{Comparison Table}

\begin{table}[h]
\centering
\caption{Comparison of Complementarity Models}
\label{tab:model-comparison}
\renewcommand{\arraystretch}{1.3}
\small
\begin{tabular}{@{}lcccccc@{}}
\toprule
\textbf{Property} & \textbf{A} & \textbf{B} & \textbf{C} & \textbf{D} & \textbf{E} & \textbf{F} \\
\midrule
Discrete domain & \checkmark & & \checkmark & & \checkmark & \checkmark \\
Continuous domain & & \checkmark & & \checkmark & & \\
Local definition & & \checkmark & & \checkmark & & \\
Global definition & \checkmark & & \checkmark & & \checkmark & \checkmark \\
Higher-order & & Limited & \checkmark & & Hypergraph & \checkmark \\
Easy estimation & & \checkmark & \checkmark & \checkmark & & \\
Parametric & & & & \checkmark & & \\
\bottomrule
\end{tabular}
\end{table}

\subsection{Equivalence Theorems}

\begin{theorem}[A $\Leftrightarrow$ B Equivalence]
\label{thm:ab-equiv}
For $f: \mathbb{R}^n \to \mathbb{R}$ with $f \in C^2$:
\begin{equation}
f \text{ is supermodular (Model A)} \Leftrightarrow \gamma^B_{ij} \geq 0 \; \forall i \neq j \text{ (Model B)}
\end{equation}
\end{theorem}

\begin{theorem}[B $\Leftrightarrow$ C Connection]
\label{thm:bc-equiv}
For small perturbations $\Delta x_i, \Delta x_j$:
\begin{equation}
\gamma^C_{ij} \approx \gamma^B_{ij} \cdot \Delta x_i \cdot \Delta x_j
\end{equation}
Factorial interaction is the discrete analogue of cross-partial.
\end{theorem}

\begin{theorem}[C $\Leftrightarrow$ F Connection]
\label{thm:cf-equiv}
Model C is Model F restricted to binary inputs:
\begin{equation}
\gamma^C_{i_1 \cdots i_k} = \Gamma_{i_1 \cdots i_k} \big|_{x_j \in \{0,1\}}
\end{equation}
\end{theorem}

\subsection{When to Use Each Model}

\begin{tcolorbox}[colback=yellow!5!white,colframe=yellow!75!black,title=\textbf{Model Selection Guide}]

\textbf{Use Model A (Supermodularity) when:}
\begin{itemize}[noitemsep]
    \item Domain has natural lattice structure
    \item Need monotone comparative statics
    \item Complementarity is global (all pairs positive)
\end{itemize}

\textbf{Use Model B (Cross-Partial) when:}
\begin{itemize}[noitemsep]
    \item Function is smooth and continuous
    \item Need local complementarity (varies by point)
    \item Doing optimization with calculus
\end{itemize}

\textbf{Use Model C (Factorial) when:}
\begin{itemize}[noitemsep]
    \item Running experiments
    \item Factors are discrete/binary
    \item Need statistical inference
\end{itemize}

\textbf{Use Model D (Production Function) when:}
\begin{itemize}[noitemsep]
    \item Building structural economic models
    \item Have theory guiding functional form
    \item Want closed-form solutions
\end{itemize}

\textbf{Use Model E (Graph) when:}
\begin{itemize}[noitemsep]
    \item Network structure is primary
    \item Analyzing social/spatial interactions
    \item Using graph algorithms
\end{itemize}

\textbf{Use Model F (Tensor) when:}
\begin{itemize}[noitemsep]
    \item Higher-order interactions matter
    \item Have high-dimensional data
    \item Need general framework
\end{itemize}

\end{tcolorbox}

% ============================================================================
% SUMMARY
% ============================================================================

%%%%%%%%%%%%%%%%%%%%%%%%%%%%%%%%%%%%%%%%%%%%%%%%%%%%%%%%%%%%%%%%%%%%%%%%%%%%%%%
\section{Key Results}
\label{sec:mdl-results}
%%%%%%%%%%%%%%%%%%%%%%%%%%%%%%%%%%%%%%%%%%%%%%%%%%%%%%%%%%%%%%%%%%%%%%%%%%%%%%%

The analysis yields the following key results:

\begin{enumerate}
    \item \textbf{Result 1:} [Primary finding with quantitative support]
    \item \textbf{Result 2:} [Secondary finding with implications]
    \item \textbf{Result 3:} [Practical application or boundary condition]
\end{enumerate}


\section{Summary}
\label{sec:summary}

\begin{tcolorbox}[colback=gray!5!white,colframe=gray!75!black,title=\textbf{Summary: Six Models of Complementarity}]

\textbf{Model A - Supermodularity:}
\begin{equation}
\gamma^A_{xy} = f(x \vee y) + f(x \wedge y) - f(x) - f(y) \geq 0
\end{equation}

\textbf{Model B - Cross-Partial:}
\begin{equation}
\gamma^B_{ij} = \frac{\partial^2 f}{\partial x_i \partial x_j}
\end{equation}

\textbf{Model C - Factorial:}
\begin{equation}
\gamma^C_{AB} = E[Y|11] - E[Y|10] - E[Y|01] + E[Y|00]
\end{equation}

\textbf{Model D - Translog:}
\begin{equation}
\ln f = \alpha_0 + \sum_i \alpha_i \ln x_i + \frac{1}{2} \sum_{ij} \gamma^D_{ij} \ln x_i \ln x_j
\end{equation}

\textbf{Model E - Graph:}
\begin{equation}
\gamma^E_{ij} = w(i, j) \text{ in graph } G = (V, E, w)
\end{equation}

\textbf{Model F - Tensor:}
\begin{equation}
\gamma^F_{i_1 \cdots i_k} = \Gamma_{i_1 \cdots i_k} \text{ in } \boldsymbol{\Gamma} \in \mathbb{R}^{n^k}
\end{equation}

\end{tcolorbox}

% ============================================================================
% GLOSSARY
% ============================================================================
\section{Glossary}
\label{sec:glossary}

\begin{description}
    \item[Supermodularity] Function property implying global complementarity
    \item[Cross-Partial Derivative] Second mixed partial derivative; local complementarity
    \item[Factorial Interaction] Deviation from additivity in experimental design
    \item[Translog] Flexible functional form with explicit complementarity parameters
    \item[Adjacency Matrix] Graph representation of pairwise complementarities
    \item[Tensor] Multi-dimensional array for higher-order complementarity
    \item[CP Decomposition] Tensor factorization into rank-one components
\end{description}

% ============================================================================
% REFERENCES
% ============================================================================

%%%%%%%%%%%%%%%%%%%%%%%%%%%%%%%%%%%%%%%%%%%%%%%%%%%%%%%%%%%%%%%%%%%%%%%%%%%%%%%
\section{Critical Foundations and Objections}
\label{sec:mdl-foundations}
%%%%%%%%%%%%%%%%%%%%%%%%%%%%%%%%%%%%%%%%%%%%%%%%%%%%%%%%%%%%%%%%%%%%%%%%%%%%%%%

\subsection{Objection 1: Scope Limitations}

\textbf{Objection:} The framework presented may have limited applicability
outside the specific domain contexts discussed.

\textbf{Response:} We acknowledge domain-specificity. The framework provides
general principles that should be calibrated to specific contexts. Cross-domain
validation remains an open research question.

\subsection{Objection 2: Measurement Challenges}

\textbf{Objection:} Key constructs may be difficult to measure reliably in
practice.

\textbf{Response:} We recommend using validated scales and instruments where
available, and developing domain-specific proxies where necessary. Sensitivity
analysis should assess robustness to measurement uncertainty.



%%%%%%%%%%%%%%%%%%%%%%%%%%%%%%%%%%%%%%%%%%%%%%%%%%%%%%%%%%%%%%%%%%%%%%%%%%%%%%%
\subsection{Core Axioms}
\label{sec:mdl-axioms}
%%%%%%%%%%%%%%%%%%%%%%%%%%%%%%%%%%%%%%%%%%%%%%%%%%%%%%%%%%%%%%%%%%%%%%%%%%%%%%%

\begin{tcolorbox}[colback=yellow!5!white,colframe=yellow!75!black,title=Axiom UNMAPPED_MDL-1: Fundamental Principle]
\textbf{Axiom UNMAPPED_MDL-1 (Core Assumption):}
The fundamental principle underlying this appendix is that behavioral outcomes
emerge from the interaction of individual characteristics and contextual factors.
\end{tcolorbox}

\begin{tcolorbox}[colback=yellow!5!white,colframe=yellow!75!black,title=Axiom UNMAPPED_MDL-2: Complementarity]
\textbf{Axiom UNMAPPED_MDL-2 (Complementarity):}
Components of the framework exhibit complementarity ($\gamma > 0$), meaning
that combined effects exceed the sum of individual effects.
\end{tcolorbox}



%%%%%%%%%%%%%%%%%%%%%%%%%%%%%%%%%%%%%%%%%%%%%%%%%%%%%%%%%%%%%%%%%%%%%%%%%%%%%%%
\section{Open Issues and Research Agenda}
\label{sec:mdl-open}
%%%%%%%%%%%%%%%%%%%%%%%%%%%%%%%%%%%%%%%%%%%%%%%%%%%%%%%%%%%%%%%%%%%%%%%%%%%%%%%

\begin{enumerate}
    \item \textbf{Empirical Validation:} Further testing across diverse contexts
    \item \textbf{Parameter Refinement:} Improved estimation methods needed
    \item \textbf{Integration:} Enhanced cross-appendix linkages
\end{enumerate}

\section{References}
\label{sec:mdl-references}

\nocite{bcm_master}

See master bibliography (\texttt{bcm\_master.bib}) for complete citations.
For comprehensive symbol definitions, see Appendix G (Master Glossary).

\subsection{Key References by Model}

\textbf{Model A:} Topkis (1998), Milgrom \& Shannon (1994), Vives (1990)

\textbf{Model B:} Samuelson (1947), Hicks (1939), Mas-Colell et al. (1995)

\textbf{Model C:} Fisher (1935), Box et al. (1978), Montgomery (2017)

\textbf{Model D:} Christensen et al. (1973), Diewert (1971), Berndt \& Christensen (1973)

\textbf{Model E:} Newman (2010), Jackson (2008), Goyal (2007)

\textbf{Model F:} Kolda \& Bader (2009), De Lathauwer et al. (2000)

% ============================================================================
% CROSS-REFERENCES
% ============================================================================
\section{Cross-References}
\label{sec:mdl-crossref}

\begin{tcolorbox}[colback=gray!5!white, colframe=gray!50!black]
\textbf{Critical Link:} \textbf{Appendix UNMAPPED_NC (FORMAL-NC)}---Non-Concavity, Fit, and Coherence.

All six canonical models share a common implication: when $\gamma > 0$, the objective function is non-concave, creating multiple local maxima. NC formalizes this connection and establishes the $K$ (coherence) vs.\ $Q$ (quality) distinction central to Chapter 7.

\textbf{Related Appendices:}
\begin{itemize}[nosep]
\item B (CORE-HOW): Complementarity $\gamma$ as the core mechanism
\item VII (FORMAL-GTC): General Theory---unified $\gamma$ framework
\item X (FORMAL-FOUND): Mathematical foundations
\end{itemize}
\end{tcolorbox}

% -----------------------------------------------------------------------------
% APPENDIX SCOPE BOX
% -----------------------------------------------------------------------------

\begin{tcolorbox}[
    colback=orange!5!white,
    colframe=orange!75!black,
    title={\textbf{Appendix Scope: Ziel / In-Scope / Out-of-Scope / Constraints}},
    fonttitle=\bfseries
]

\textbf{Ziel (Objective):}
\begin{itemize}[nosep]
    \item Formalize and document the key concepts of Mathematical Models of Complementarity: A Comprehensive Catalog
\end{itemize}

\textbf{In-Scope:}
\begin{itemize}[nosep]
    \item Core theoretical foundations
    \item Mathematical formalizations where applicable
    \item Integration with EBF framework
\end{itemize}

\textbf{Out-of-Scope (delegiert an andere Appendices):}
\begin{itemize}[nosep]
    \item Detailed empirical validation $\rightarrow$ METHOD appendices
    \item Domain-specific applications $\rightarrow$ DOMAIN appendices
    \item Complete literature review $\rightarrow$ LIT appendices
\end{itemize}

\textbf{Constraints (Anwendungsgrenzen):}
\begin{itemize}[nosep]
    \item Framework requires calibration to specific contexts
    \item Parameters may vary across domains
    \item Validity bounded by underlying assumptions
\end{itemize}

\textbf{Lieferobjekte (Deliverables):}
\begin{itemize}[nosep]
    \item Theoretical framework with formal definitions
    \item Integration points with other appendices
    \item Practical guidelines for application
\end{itemize}

\end{tcolorbox}

%%%%%%%%%%%%%%%%%%%%%%%%%%%%%%%%%%%%%%%%%%%%%%%%%%%%%%%%%%%%%%%%%%%%%%%%%%%%%%%
\section{The Fundamental Question}
\label{sec:mdl-fundamental}
%%%%%%%%%%%%%%%%%%%%%%%%%%%%%%%%%%%%%%%%%%%%%%%%%%%%%%%%%%%%%%%%%%%%%%%%%%%%%%%

\begin{tcolorbox}[
    colback=red!5!white,
    colframe=red!75!black,
    title={\textbf{Fundamental Question}}
]
\textbf{How does mathematical models of complementarity: a comprehensive catalog integrate with and contribute to the
Evidence-Based Framework for understanding behavioral outcomes?}

This appendix addresses the core challenge of formalizing behavioral principles
within a rigorous theoretical framework while maintaining practical applicability.
\end{tcolorbox}





\end{chapter}
